% preamble:
\documentclass[11pt,a4paper, titlepage]{article}
\usepackage[a4paper, margin=0.85in]{geometry}  % Adjust margin size
% \usepackage{amsmath}
% \usepackage{graphicx}
% \usepackage{hyperref}


\author{Davide Squeri}
\title{Documentazione Progetto Reti Informatiche\\[0.3em]
\large{Anno Accademico 2025/2026, matricola pari}}
\date{19/02/2025}

% document
\begin{document}
\maketitle

\section*{Documentazione}
\subsection*{Scelta del Protocollo di Trasporto}
TCP per le comunicazioni Utente - Lavagna. La TCP handshake funziona da HELLO, o meglio il comando HELLO inizia la TCP handshake.
UDP per le comunicazioni Utente - Utente, poichè le comunicazioni tra utenti sono sporadiche, così da risparmiare l'overhead della connessione TCP.
\subsection*{Scelta del tipo di concorrenza per la Lavagna}
Ho optato per rendere la lavagna iterativa (non concorrente) nel servire le connessioni con gli utenti poiché:
\begin{enumerate}
    \item l'approccio iterativo risparmia overhead
    \item la sincronizzazione è più semplice
    \item non serve concorrenza perché la lavagna esegue solo operazioni I/O bound leggere
\end{enumerate}
Tengo traccia degli utenti con una lista senza limite massimo; con un solo thread multiplexo con \texttt{select()} stdin, nuove connessioni ed esistenti, servendoli in quest'ordine. Ogni scambio utente-lavagna deve quindi essere botta-e-risposta senza input intermedio, altrimenti la lavagna resterebbe bloccata su un utente: per esempio l'esito della \texttt{CREATE\_CARD} (ID duplicato o no) arriva solo dopo l'invio completo, e scomporlo in pi\`u fasi aumenterebbe le \texttt{send}/\texttt{recv}.

\subsection*{Struttura dei pacchetti e modalità di scambio dei Dati}
Per quanto riguarda le comunicazioni utente - lavagna, avendo optato per chiudere gli scambi di informazioni in al più un messaggio da parte dell'utente uno da parte della lavagna, ho la necessità di inviare più informazioni, anche di natura diversa in una stessa \texttt{send()}.

Ho deciso che i primi 4 byte di ogni messaggio sono l'ID del comando (\texttt{uint32\_t}), in base a quello la lavagna (o l'utente) sa come interpretare il resto del messaggio. Per esempio la \texttt{CREATE\_CARD} invia un messaggio costruito come segue:
\begin{verbatim}
4byte: ID_CREATE_CARD
4byte: len_descrizione (strlen+1)
4byte: ID_nuova_card
4byte: colonna (0=ToDo, 1=Doing, 2=Done)
max1024byte: descrizione_attività
1byte:'\0'
\end{verbatim}
Il campo lunghezza permette al ricevente di sapere esattamente quanti byte leggere (letture in ciclo fino a esaurimento). La \texttt{HELLO} \`e \texttt{[ID\_HELLO][porta\_utente]} con risposta \texttt{"OK"}/\texttt{"NO"}; la \texttt{QUIT} \`e il solo ID comando. La \texttt{AVAILABLE\_CARD} arriva dopo una \texttt{SEND\_USER\_LIST} per-recipient (lista degli altri idonei + conteggio), poi \texttt{[ID\_AVAILABLE\_CARD][len][ID\_card][descrizione]}. L'asta \texttt{CHOOSE\_USER} viaggia su UDP con pacchetti \texttt{[ID\_card][porta][costo]}: l'ID card scarta i pacchetti stantii delle aste precedenti.

I numeri sono portati in endianess di rete prima di essere inviati, e riportati in endianess di sistema all'arrivo.

\subsection*{Gestione dei socket}
La lavagna usa un socket TCP di ascolto e uno slot per utente, tutti multiplexati con \texttt{select()} a timeout di 1s assieme a stdin: accetta connessioni, serve tutti i descrittori pronti (senza affamarne alcuno) ed esegue il controllo periodico dei PING. Ogni utente ha un socket TCP verso la lavagna e un socket UDP verso i pari per l'asta.

\subsection*{PING\_USER e PONG\_LAVAGNA}
Le card in \texttt{Doing} da pi\`u di 90s ricevono \texttt{[ID\_PING\_USER][ID\_card]} e l'utente risponde \texttt{[ID\_PONG\_LAVAGNA][ID\_card]}; senza PONG entro 30s la card torna in \texttt{ToDo} (soglie in \texttt{include.h}). Pregio: nessun thread dedicato, solo il ciclo \texttt{select()}. Difetto: controllo O(n) a ogni ciclo e risposta differita se l'utente \`e in asta (comunque entro i 30s).

\section*{Scelte Varie}
Ho optato per non rendere bloccante l'assegnazione di una Card: l'asta UDP ha timeout di 10s e ignora i pacchetti delle aste precedenti, cos\`i un utente che esce a met\`a asta non blocca gli altri.

All'avvio la lavagna precarica 10 card in \texttt{ToDo} come richiesto dalle specifiche; la \texttt{QUIT} di un utente rimette le sue card \texttt{Doing} in \texttt{ToDo} per riassegnarle.

Se all'interno di un'asta, un utente smette di rispondere, viene assunto disconnesso dagli altri utenti e l'asta continua.

Utenti già impegnati con una card in Doing non partecipano alle aste.

Ho scelto come range di tempo necessario per card [1 ; 100] secondi, possono sembrare lunghi ma serve che qualche card richieda più di 90 secondi altrimenti non è verificabile il funzionamento di PING e PONG.

\section*{Comandi disponibili}
Avvio: \texttt{./lavagna} (nessun argomento, porta 5678) e \texttt{./utente <porta>} (porta $\geq$ 5679).
\subsection*{Lavagna (da terminale)}
\begin{itemize}\setlength{\itemsep}{1pt}\setlength{\parsep}{0pt}
    \item \texttt{SHOW\_LAVAGNA}: ristampa la lavagna con colonne e dettaglio card.
    \item \texttt{QUIT}: termina la lavagna.
\end{itemize}
Gli altri comandi digitati sulla lavagna rispondono ``non disponibile'': \texttt{ACK\_CARD}, \texttt{CARD\_DONE} e \texttt{MOVE\_CARD} avvengono in automatico alla ricezione dei messaggi degli utenti.
\subsection*{Utente (da terminale)}
\begin{itemize}\setlength{\itemsep}{1pt}\setlength{\parsep}{0pt}
    \item \texttt{HELLO}: prima operazione obbligatoria, connette e registra l'utente (una sola volta).
    \item \texttt{CREATE\_CARD}: chiede ID, colonna (0=ToDo, 1=Doing, 2=Done) e descrizione; riporta accettata o ID duplicato.
    \item \texttt{QUIT}: notifica l'uscita alla lavagna e termina (le sue card \texttt{Doing} tornano in \texttt{ToDo}).
\end{itemize}
\texttt{ACK\_CARD} e \texttt{CARD\_DONE} sono automatici (vincita asta e fine lavoro); \texttt{SEND\_USER\_LIST}, \texttt{AVAILABLE\_CARD}, \texttt{PING\_USER} e \texttt{PONG\_LAVAGNA} sono solo protocollo e digitati rispondono ``non disponibile''.

\section*{Istruzioni sulla Compilazione}
Comando per compilare il progetto:
\begin{verbatim}
make all
\end{verbatim}
Comando per rimuovere i file oggetto:
\begin{verbatim}
make clean
\end{verbatim}
Il comando crea i due file eseguibili richiesti nella root directory del progetto.
Per eseguirli si può proseguire come indicato nelle specifiche del progetto.


\end{document}
